\section{价格的延拓与同时极小化}\label{sec:prices}

压缩论证需要引入新站点来安放张量的腿。旧站点上的价格已经为任一拟议放置方式的价格给出了下界。具体地说，在新站点收缩各腿，再对剩余张量求值。对任意张量幂都可如此操作，每个新站点可以一次性收缩它的所有副本。定理~\ref{thm:extension} 表明这些下界是精确的，并给出了同一放置方式同时满足多个价格界的判据。在第~\ref{sec:sync} 节，待放置的张量将依赖于一个参数。为了控制该参数的特殊化，我们在赋值域上证明这一判据。对于与一个新站点关联的若干普通连线的乘积，其最小值已在 \cite[Theorem~\ext{G}{thm:extension}]{oaG} 中算出。

\subsection{整限制}
设 \(K\) 是带实值非阿基米德加法赋值 \(\nu\) 的域，约定 \(\nu(0)=+\infty\)。其赋值环为 \(\mathcal O=\{c\in K:\nu(c)\ge0\}\)。各站点均采用给定的坐标基。如果每个局部矩阵的元素均在 \(\mathcal O\) 中，就称该限制为\emph{整限制}。定义在 \(S\) 上、尺度为 \(D>0\) 的\emph{赋值价格} \(p\)（仍假设 $|S|\ge2$）在非零张量上取实值，满足 \(p(0)=-\infty,\ p(\mathbf1)=0\)，且在整限制下不增。其标量、乘积与直和规则为
\begin{equation}\label{eq:prices-axioms}
 \begin{gathered}
 p(cT)=p(T)-\nu(c) \quad(c\ne0), \qquad
 p(T\otimes U)=p(T)+p(U),\\
 p\Bigl(\bigoplus_j T_j\Bigr)=D\log\sum_j e^{p(T_j)/D}.
 \end{gathered}
\end{equation}
这里 \(\log\) 表示自然对数。因此乘以赋值为正的标量会降低价格。取平凡赋值即恢复普通价格。

我们将利用直和规则的一种方便的等价形式来构造价格。对于处在相同空间中的张量，逐站点将各块相加，便知 \eqref{eq:prices-axioms} 右端给出了 \(p(\sum_j T_j)\) 的上界。反过来，假设此普通和的上界成立，并且对每个 \(q\) 都有 \(p(\Delta_S(q))=D\log q\)。通过各块的嵌入可将 \(\bigoplus_j T_j\) 写成普通和，从而得到所需上界。为证下界，取直和的 \(m\) 次幂，在每个站点投影到这样一些词上：第 \(j\) 个直和项在其中出现 \(m r_j\) 次。经局部置换，所得张量可认作 \(\bigotimes_jT_j^{\otimes m r_j}\) 与大小为 \(m!/\prod_j(mr_j)!\) 的对角张量的乘积。对于支撑在非零直和项上的有理概率向量 \(r\)，让 \(m\) 取其公分母的各倍数。每份副本的价格至少为 \(\sum_jr_jp(T_j)-D\sum_j r_j\log r_j\)。取最大值即得所述等式。

以 \(q_i\) 表示 \(T\) 的局部维数，以 \(\nu_{\min}(T)\) 表示其各分量赋值的最小值。当 \(T\ne0\) 时，
\begin{equation}\label{eq:prices-range}
 -\nu_{\min}(T)\le p(T)\le -\nu_{\min}(T)+D\sum_i\log q_i.
\end{equation}
投影到一个具有最小赋值的分量即得下界。将 \(T\) 展开成至多 \(\prod_iq_i\) 个坐标张量并应用普通和的不等式即得上界。由此，固定尺度的价格在逐点收敛拓扑下构成紧空间，因为其公理在区间族 \eqref{eq:prices-range} 的乘积中给出闭条件。前述连线规则和搬运规则只需整运算，故对 \(p\) 同样成立。在单个站点，非零向量的价格是 \(-\nu_{\min}\)，因为除以一个具有最小赋值的分量后，它就与单位张量整等价。当 \(K\) 包含一个赋值平凡的子域 \(F\) 时，限制到 \(F\) 上的张量就得到连线距离的普通谱几何。

\subsection{价格延拓的判据}
固定 \(S\) 上的 \(p\)，并添加由有限个站点组成的集合 \(L\)。以 \(\widehat p\) 表示到 \(S\sqcup L\) 的同尺度延拓，要求它在 \(L\) 处仅有平凡腿的张量上保持 \(p\)。将新增站点放在旧站点上即得这样的延拓。限制到任何包含 \(S\) 的中间站点列表，所得都是尺度为 \(D\) 的赋值价格。事实上，该列表上的对角张量可限制为基底的对角张量，它本身又是完整对角张量的限制，因此可应用对角张量判据。

一次\emph{闭合}在 \(L\) 的每个站点上施以一个整线性泛函，并保留基底的所有自由腿。\(T\) 的\emph{辅助成本}为
\begin{equation}\label{eq:prices-beta}
 \beta_p(T)=\sup_{m\ge1}\frac1m
       \sup_{\ell}p\bigl(\ell(T^{\otimes m})\bigr).
\end{equation}
在维数为 \(q\) 的站点，上式中的泛函作用于 \(K^{q^m}\)，因而可任意耦合这 \(m\) 个副本。由闭合下的单调性与张量幂上的可加性，每个延拓都满足 \(\widehat p(T)\ge\beta_p(T)\)。

例如在普通情形下，考虑两条大小均为 \(q\) 的连线的乘积，其中一条从 \(a\) 连到 \(y\)，另一条从 \(y\) 连到 \(b\)，\(a,b\) 是旧站点，\(y\) 是新站点。只需一个副本，利用配对 \(e_i\otimes e_j\mapsto\delta_{ij}\) 闭合 \(y\) 处的两条腿，就可恢复从 \(a\) 到 \(b\) 的连线。这给出乘积的下界 \(d(a,b)\log q\)。将 \(y\) 放在 \(a\) 即可达到此界。更一般地，沿测地线段距离可加，故将 \(y\) 放在从 \(a\) 到 \(b\) 的测地线上任一点都可达到此界。

\begin{theorem}[价格延拓]\label{thm:extension}
对 \(S\sqcup L\) 上的任意张量 \(T\)，有 \(\beta_p(T)=\min_{\widehat p}\widehat p(T)\)。给定任意一族非零张量 \(T_i\) 和实数 \(h_i\)，存在使得对所有 \(i\) 均有 \(\widehat p(T_i)\le h_i\) 的同一个延拓，当且仅当对每个有限指标列表（允许重复）都有
\begin{equation}\label{eq:prices-product-test}
 \beta_p(T_{i_1}\otimes\cdots\otimes T_{i_j})
 \le h_{i_1}+\cdots+h_{i_j}
\end{equation}
\end{theorem}

\begin{proof}
\emph{单个张量的极小化。} 零张量的情形显然。给定非零张量 \(T\)，我们先向包含 \(T\) 的、由互不相同的非零张量组成的有限列表 \(A_1,\ldots,A_s\) 赋予价格。以它们的混合乘积中的指数为变量。对 \(x\in\mathbb Q_{\ge0}^s\)，令
\[
 b(x)=\frac1n\beta_p\Bigl(\bigotimes_i A_i^{\otimes nx_i}\Bigr),
 \qquad nx\in\mathbb Z^s.
\]
在 \eqref{eq:prices-beta} 中，独立重复每一次闭合，便知可把上确界限制在任一给定正整数的所有倍数上。因此辅助成本对张量幂是齐次的，从而 \(b\) 定义良好。在共同的倍数处取独立闭合可知 \(b\) 具有超可加性，再由齐次性可知它是凹函数。

坐标 \(i\) 每增加一单位所产生的增量落在 \(A_i\) 的价格区间 \eqref{eq:prices-range} 内。为得下界，在基底站点和新增站点上作坐标投影，从每个新增因子中取出一个赋值最小的系数。为得上界，将每个新增因子按坐标展开。即使闭合耦合了所有副本，固定各项的新增指标后，仍在旧因子上诱导出一个整闭合。基底站点上留下的坐标因子价格为零，由普通和的上界，每新增一个形状为 \((q_j)\) 的因子，所付代价至多为 \(D\sum_j\log q_j\)。取张量幂即得所述增量界。特别地，\(b\) 可连续延拓到实数非负正交象限。

在所有坐标严格为正的有理点 \(x\)，取这个齐次凹函数的一个支撑线性形式 \(u\)，使得 \(b(y)\le u\cdot y\) 且 \(b(x)=u\cdot x\)。双侧增量估计表明 \(u_i\) 落在同一价格区间中。我们断言，各数 \(u_i\) 满足列表上的价格关系。例如，若 \(A_i=A_j\otimes A_k\)，则对绝对值充分小的正、负有理数 \(t\) 都有 \(b(x+t(e_j+e_k-e_i))=b(x)\)。沿这两个方向的支撑不等式迫使 \(u_i=u_j+u_k\)。作标量变换时，\(b\) 对每份副本的改变量为该标量赋值的相反数，同理即得标量关系。若 \(A_i\) 经整限制可得 \(A_j\)，则将若干 \(A_j\) 的副本换成 \(A_i\) 的副本，只会使 \(b\) 增大或不变。故沿 \(e_i-e_j\) 方向的支撑不等式给出 \(u_i\ge u_j\)。每添加一个单位的基底张量，\(b\) 恰好增加该张量的价格，因为其因子在每次闭合中均原样通过。对于 \(\Delta_{S\sqcup L}(q)\)，在其新增腿上用全一泛函闭合便得 \(\Delta_S(q)\)，而将其 \(q\) 个对角分量展开即得反向不等式。因此其增量为 \(D\log q\)。

为验证普通和的上界，设列表中有 \(A_i=A_j+A_k\)。取有理数 \(0<\sigma\le x_i\)，在重数为 \(nx\) 的乘积中展开 \(A_i\) 的 \(n\sigma\) 个副本。选取 \(A_j\) 的 \(a\) 个副本的项，闭合后的价格至多为 \(n u\cdot x-n\sigma u_i+a u_j+(n\sigma-a)u_k\)。利用基底上普通和的上界以及二项式公式，得到原张量闭合后的价格至多为
\[
 n u\cdot x-n\sigma u_i+
 n\sigma D\log\bigl(e^{u_j/D}+e^{u_k/D}\bigr).
\]
令 \(n\) 取各个允许的倍数，将价格除以它后取上确界，左侧即为 \(b(x)=u\cdot x\)。这证明了关于 \(u_i\) 的所需上界。含有零加数的关系自动成立。

令 \(x\) 趋于 \(T\) 所对应的单位向量。支撑向量有收敛子列，由连续性，极限满足 \(u_T=\beta_p(T)\)。因此，对任意有限组价格关系及 \(T\) 处的等式，只要对这些约束所涉及的全部张量组成的列表应用上述构造，就得到一个可行解。由区间乘积的紧性，所有这些关系有公共解。利用 \eqref{eq:prices-axioms} 后所述的对角张量判据，此解便是赋值价格。这就证明单个张量的最小值可以达到。

\smallskip\noindent
\emph{同时满足多个界。} 乘积上的可加性表明，任何满足所给上界的延拓都满足 \eqref{eq:prices-product-test}。反之，取一张含 \(k\) 个界的有限列表。对每个 \(i\) 取 \(T_i^{\otimes n}\) 的 \(N_i\) 个副本，再将所有副本作直和，记为 \(U\)。在 \(U\) 的张量幂中，由各直和项组成的每个词都是一个混合乘积。将一个闭合限制在该词所对应的各块上，便得整闭合，它的价格受 \eqref{eq:prices-product-test} 控制。在基底上对各词求和，并按每个副本取上确界，得到
\[
 \beta_p(U)\le D\log\sum_i N_i e^{nh_i/D}.
\]
取 \(C>\max_i h_i\) 及 \(N_i=\lceil e^{n(C-h_i)/D}\rceil\)。这样，各项 \(N_i e^{nh_i/D}\) 都属于 \([e^{nC/D},2e^{nC/D}]\) 内。对于使 \(U\) 取最小值的延拓，直和等式随即给出
\[
 D\log N_i+n\widehat p(T_i)\le\beta_p(U)\le nC+D\log(2k),
 \qquad
 \widehat p(T_i)\le h_i+\frac{D\log(2k)}{n}.
\]
令 \(n\) 趋于无穷，先对有限列表利用紧性，再通过有限交性质推广至全族。
\end{proof}

下文将用到两个推论。给定 \(p\) 的若干延拓并固定其任一凸组合，则存在同一个延拓，使其在每个张量上的取值均不超过这些延拓取值的该凸组合。事实上，各延拓都给出每个混合乘积的辅助成本上界，故将各界加权平均即可通过判据 \eqref{eq:prices-product-test}。也可在应用判据之前扩大基底。假设扩大后的基底带有一个在旧基底上与 \(p\) 一致的价格。支撑在旧基底和未知站点上的张量，在额外的基底站点仅有平凡腿，故其闭合与此前完全相同，从而其辅助成本不变。当扩大后的基底尺度更大时，这一观察仍成立。于是定理~\ref{thm:extension} 给出该较大尺度下的延拓。

\subsection{沿测地线放置}
本节以下假设 \(K\) 包含赋值平凡的子域 \(F\)。给定若干站点对 \(a_j,b_j\)，先把新站点 \(y_j\) 分别放在 \(a_j\) 处，再分别放在 \(b_j\) 处。定理~\ref{thm:extension} 的插值推论给出一个延拓，使每个张量的价格不超过上述两次求值以 \(1-\theta,\theta\) 为权重所作的仿射插值。特别地，
\[
 d(a_j,y_j)\le\theta\,d(a_j,b_j),\qquad
 d(y_j,b_j)\le(1-\theta)d(a_j,b_j).
\]
三角不等式迫使两个界都取等号。因此，在 \(F\) 上的一个完备 \(\operatorname{CAT}(0)\) 谱实现中，\(y_j\) 是沿从 \(a_j\) 到 \(b_j\) 的测地线走过全程 \(\theta\) 处的点。同一延拓给出了 \(K\) 上所有张量的插值上界。

\begin{lemma}[测地决定性与次平均不等式]\label{lem:geodesic}
完备 \(\operatorname{CAT}(0)\) 谱空间中任意子集上的普通价格数据决定其闭测地凸包上的所有价格，至多相差一个固定给定数据的保价格等距映射。此外，假设对于原始张量槽位上的每个张量，中心放置下的价格均不超过有限多个端部放置下价格的加权平均。在每种放置中，往对应的测地线段上添加点。如果在中心放置中，各新点到其两个端点的距离分别是各端部放置下相应距离的加权平均，则扩大后的各放置仍满足该价格不等式。
\end{lemma}

\begin{proof}
对于每个新站点，指定它到两个端点的大小为二的连线的价格。这些价格等于两段测地距离各自乘以 \(\log 2\)。我们断言，扩大后的列表上任一 \(T\ne0\) 的价格，等于使上界 \(h\) 与所有连线的上界满足乘积判据的最小数 \(h\)。由闭合可知实际价格满足判据。反之，扩大基底以纳入实际的测地线上的点，对具有上述上界的副本站点应用定理~\ref{thm:extension}。在共同的度量实现中，三角不等式迫使每个副本到两个端点的距离均等于指定距离。由测地线的唯一性，该副本即为实际的点。于是由搬运关系 \eqref{eq:prices-transport} 得到断言。对坍缩的基底，其所有指定线段均为常值，断言直接成立。

迭代使用这一刻画，即可确定所有有限次测地构造上的价格，包括距离。由测地插值的连续性，这些构造的并集的完备化就是闭凸包；利用 \eqref{eq:prices-transport} 又可将价格的一致性推广到完备化。原子集无限也无妨，因为每次有限构造都从有限多个点出发。

对于第二个断言，中心放置坍缩的情形显然。否则，固定判据中的一个乘积、它的一个张量幂及其新站点上的一次闭合。由原次平均不等式，闭合张量在中心的价格不超过它在各端部闭合后价格的加权平均。在每个端部，由限制下的单调性，闭合后的价格不超过闭合前该张量幂的价格；这也包括坍缩的端部，因为此时每个非零张量的价格均为零。除以幂次并在中心取上确界。因此，在扩大后的放置上，将实际的上界加权平均，就能在中心通过乘积判据，而加权平均后的连线上界又等于中心指定的界。应用刚证明的刻画即得结论。
\end{proof}

\subsection{极小化站点处求值的唯一性}
如两条连线的例子所示，价格的最小值可以在多种放置下达到。为了比较参数变化时的极小化放置，我们在每个未知站点引入一个严格凸的惩罚项。在基底上固定 \(o\)。对于新站点 \(y\)，取 \(\mathbf1\) 与从 \(o\) 到 \(y\) 的大小为二的连线的直和。其直和标号在配置的每个站点上均可见，因此其价格是
\[
 f(r)=D\log\bigl(1+\exp(r\log(2)/D)\bigr),\qquad r=d(o,y).
\]
对于基底和一列新站点上的非零张量 \(U\)，固定整数 \(M\ge1\)。其\emph{钉扎目标} \(V\) 是 \(U^{\otimes M}\) 与每个新站点各自的一个此类惩罚张量的乘积。如有多个目标，则每个目标都在基底和自身的站点列表上构造，再在其余列表的站点处补上平凡腿。

\begin{lemma}[钉扎求值]\label{lem:pinning}
对于定义在两两不交的新站点列表上的有限族目标，存在同时使各目标达到最小值的同尺度延拓。所有这样的延拓对合并列表上的每个张量均取相同的值。这些值在各子族的重叠部分一致。对于任意这样的族，包括无限族在内，所得的对 \(F\) 上张量的求值可在一个完备的 \(\operatorname{CAT}(0)\) 谱空间中实现。

令 \(W\) 为某个有限族目标的乘积，令 \(G=\sum_V\beta_p(V)\)，并令 \(h_T\) 为非零测试张量 \(T\) 的钉扎求值。存在仅依赖于 \(D\) 的 \(c_D>0\)，使得对每个 \(n\ge1\) 都有
\begin{equation}\label{eq:prices-pin-limit}
 0\le h_T-\bigl(\beta_p(W^{\otimes n}\otimes T)-nG\bigr)
 \le\frac{\sum_j(\log q_j)^2}{8c_D n},
\end{equation}
其中 \(q_j\) 是在各新增站点将测试腿合并后的维数；未用到的站点计为一维。
\end{lemma}

\begin{proof}
依次扩大基底，并用定理~\ref{thm:extension} 使下一个目标极小化。由扩大基底的性质，每次都能达到原有的最小值。

为了在一个共同的延拓中比较两种放置，取所有对应站点的中点作插值，使 \(U\) 的赋值价格不超过它在两种放置下价格的平均值。对任意两站点 \(a,b\) 及其中点 \(m\)，有
\begin{equation}\label{eq:prices-midpoint}
 f(d(o,m))\le \tfrac12\bigl(f(d(o,a))+f(d(o,b))\bigr)
                  -c_Dd(a,b)^2 .
\end{equation}
事实上，利用完整对角张量的限制可知，连线距离至多为 \(D\)。利用 \(f'\) 和 \(f''\) 在 \([0,D]\) 上的正下界，取足够小的 \(c_D>0\)，使 \(f(r)-4c_Dr^2\) 在该区间上凸且单调不减。对该函数应用中点凸性，对 \(4c_Dr^2\) 应用距离平方的中点不等式。每个目标的站点均包含原基底，因此即使在更大的配置中，其直和惩罚项的尺度仍是 \(D\)。如果两种放置均使此目标取到最小值，则中点的值不超过同一最小值减去 \(c_D\) 乘以所有对应新站点间距离的平方和。因此这些距离必定全部为零。

给定两个极小化延拓 \(P,P'\)，向 \(P\) 添加整个列表的第二份副本，并用 \(P'\) 对其求值施加上界。这些界在旧基底上通过乘积判据，故在新基底上同样通过。每个复制的目标均达到其最小值，因此它的各站点与对应的原站点之间距离为零。由搬运关系，对于每个测试张量有 \(P(T)\le P'(T)\)。交换 \(P,P'\) 即得等号。限制到任一子族后，其中的各目标仍均达到最小值。所得各有限列表上的普通求值彼此相容，按零距离取商即给出一个谱空间，从而得到所述的完备延拓。

最后，如上向任意一个参与比较的放置添加一个最优副本，并令 \(d_j\) 表示对应站点之间的距离。中点的目标值至少是 \(G\)，故由 \eqref{eq:prices-midpoint} 得 \(p_{\rm comp}(W)-G\ge 2c_D\sum_jd_j^2\)。搬运关系给出 \(p_{\rm comp}(T)\ge h_T-\sum_jd_j\log q_j\)。将第一个不等式乘以 \(n\)，加上第二个不等式，再对各 \(d_j\) 配方。利用定理~\ref{thm:extension} 取最小值，即证得 \eqref{eq:prices-pin-limit} 的上界。最优放置本身给出下界。
\end{proof}
